%-*-latex-*-
%-*-latex-*-
\newcommand\COURSE{ciss245}
\newcommand\ASSESSMENT{q4101}
\newcommand\ASSESSMENTTYPE{Quiz}
\newcommand\POINTS{\textwhite{xxx/xxx}}

\input{myquizpreamble}
\input{yliow}
\input{\COURSE}
%-*-latex-*-

%-*-latex-*-

%-*-latex-*-
\renewcommand\TITLE{\ASSESSMENTTYPE \ \ASSESSMENT}

\renewcommand\EMAIL{}
\renewcommand\AUTHOR{eagraham1@cougars.ccis.edu}

\textwidth=6in
\begin{document}
\topmatterthree

%------------------------------------------------------------------------------
\nextq
Write a struct template that models polynomials.
A polynomial is an array of coefficients.
For instance $1 + 2x + 42x^2$ can be modeled with ${1, 2, -42}$.
\\
\ANSWER
\begin{answercode}
#include <iostream>

template < typename T >
struct polynomial
{
    T coef[1024];
    int size;
};

template < typename T >
void init(polynomial< T > & p, T x[], int size)
{
  p.size = size;
  for (int i = 0; i < size; ++i)
  {
    p.coef[i] = x[i];
  }
}

template < typename T >
void println(polynomial< T > & p)
{
  bool first = true;
  for (int i = 0; i < p.size; ++i)
  {
    if (p.coef[i] != 0)
    {
      if (!first)
      {
        std::cout << " + ";
      }
      std::cout << p.coef[i];
      if (i == 1)
      {
        std::cout << "x";
      }
      else if (i > 1)
      {
        std::cout << "x^" << i;
      }
      first = false;
    }
  }
  if (first)
  {
    std::cout << "0";
  }
  std::cout << std::endl;
}

int main()
{
    polynomial< int > p;
    int x[] = {1, 2, 42};
    init(p, x, 3);
    println(p); // prints 1 + 2x + 42x^2 (and newline)

    polynomial< double > q;
    double y[] = {1.1, 0.0, 42.42, 0.0, 0.0, 0.0, 0.0};
    init(q, y, 7);
    println(q); // prints 1.1 + 42.42x^2 (and newline)

    return 0;
}
\end{answercode}

%------------------------------------------------------------------------------
%-*-latex-*-
\newpage
\input{instructions}
\end{document}

